\chapter{Kết luận}
\label{Chapter6}

\textit{Chương 6 sẽ trình bày những kết quả đạt được sau khi hoàn thành xong đề tài đồ án này, định hướng phát triển của ứng dụng trong tương lai.}

\section{Kết quả đạt được}

Xuyên suốt quá trình thực hiện đồ án tốt nghiệp trong 6 tháng qua, nhóm đã học tập và tích luỹ được rất nhiều kiến thức tuy rằng không được đào tạo trong chuyên ngành nhưng rất bổ ích cho việc trau dồi kiến thức về ngành Công Nghệ Thông Tin.

\subsection{Về mặt kiến thức}
\begin{itemize}
	\item [-] Tìm hiểu và áp dụng thành công kĩ thuật Scraping Web để thu thập dữ liệu từ các web site.
	\item [-] Nghiên cứu các thư viện tách từ theo các từ ghép như underthesea, VnCoreNLP, python vietnamese toolkit.
	\item [-] Tìm hiểu và áp dụng được pre-trained Pho BERT, ngoài ra còn tự build được một mô hình của riêng nhóm trên dữ liệu tự tạo.
	\item [-] Sử dụng được Elastic Search để xử lý chính cho hệ thống.
	\item [-] Xây dựng được Web Server cho hệ thống.
	\item [-] Giải quyết được tất cả các bài toán cần phải giải quyết để xây dựng hệ thống.
\end{itemize}

\subsection{Về mặt ứng dụng}

Nhóm đã có thử đưa ứng dụng lên một máy chủ VPS Azure để có thể host ứng dụng thành một ứng dụng Web hoàn chỉnh nhưng vì lý do chi phí để mở máy chủ quá cao và chưa thành công trong việc đăng kí tên miền cho máy chủ nên hiện về mặt ứng dụng thì vẫn chưa thành công. Nhóm có thử sử dụng qua các dịch vụ miễn phí hoặc là có ưu đãi cho sinh viên nhưng vẫn không thể đáp ứng được việc đủ bộ nhớ để lưu dữ liệu Elasticsearch và các thư viện của hệ thống, bộ nhớ Ram phục vụ cho Elasticsearch hoạt động.
\section{Định hướng phát triển}

Vì giới hạn về kĩ năng cũng như thời gian, hiện tại nhóm đã hoàn thành tốt nhất sản phẩm cho hiện tại nhưng vẫn còn một số hạn chế. Trong tương lai, nhóm sẽ phát triển một số tính năng sau nữa để có thể hoàn thành nhất sản phẩm ví dụ như:
\begin{itemize}
	\item [-] Xây dựng thêm tính năng thêm số lượng câu hỏi và câu trả lời vào hệ thống, vì hiện tại hệ thống chỉ sử dụng trên dữ liệu được lưu trữ sẵn trong cơ sở dữ liệu và không thể thêm dữ liệu mới vào vậy nên khi cần thêm mới bộ câu hỏi và câu trả lời thì nhóm cần phải xây dựng thêm tính năng mới để có thể thêm mới dữ liệu vào được.
	\item [-] Nâng cấp được độ chính xác của mô hình đào tạo trên dữ liệu Pháp luật. Hiện tại thì mô hình vẫn chỉ dừng lại ở mức độ là kiểm tra câu hỏi của người dùng có tương đồng với dữ liệu có trong cơ sở dữ liệu hay không chứ chưa thể kiểm tra được câu hỏi của người dùng đồng nghĩa hay trái nghĩa với câu hỏi được lưu trữ trong cơ sở dữ liệu. Vậy nên, định hướng để đưa mô hình đạo tạo chính xác hơn thì nhóm cần phải thực hiện được chức năng này.
	\item [-] Triển khai được hệ thống hoạt động như Web app hoàn chỉnh. Hiện tại vì lý do kinh phí duy trì được lượng dữ liệu trong cơ sở dữ liệu và sử dụng Elasticsearch để tìm kiếm cũng như có đủ phần cứng để chạy được mô hình pre-trained PhoBERT thì cần phải có một phần cứng đủ mạng để có thể chịu được hệ thống. Nhưng do kinh phí để có thể duy trì được một hệ thống như vậy tương đối lớn so với nhóm và cũng như nhóm thể làm kịp được việc triển khai hệ thống nên phần này nhóm sẽ định hướng để phát triển thêm.
\end{itemize}

